\history{Manuscript revised September 13, 2026.}
\doi{}

\title{PGA-UNet: A Lightweight Prompt-Guided Attention U-Net for Interactive Bone X-Ray Lesion Segmentation}
\author{\uppercase{Thong Luc}\authorrefmark{1,2},
\uppercase{Huu-Binh Nguyen}\authorrefmark{1,2},
\uppercase{Ly Quoc Ngoc}\authorrefmark{1,2}}

\address[1]{B.Sc. Program in Computer Science, Faculty of Information Technology, University of Science, VNU-HCM, Vietnam}
\address[2]{Viet Nam National University, Ho Chi Minh City, Vietnam}

\markboth
{Luc \headeretal: PGA-UNet for Interactive Bone X-Ray Lesion Segmentation}
{Luc \headeretal: PGA-UNet for Interactive Bone X-Ray Lesion Segmentation}

\corresp{Corresponding authors: Ly Quoc Ngoc (e-mail: lqngoc@fit.hcmus.edu.vn) and Thong Luc (e-mail: 22120196@hcmus.edu.vn).}

\begin{abstract}
Segmenting bone lesions in X-ray images can be difficult when lesions occupy a small image area or have low contrast. We present PGA-UNet, a compact convolutional network for segmentation from an image and a lesion-covering box. A Gaussian-smoothed plateau encodes the box and conditions encoder features through a Prompt Spatial Gate and decoder skip attention through a Conditional Attention Decoder. We evaluate separate models on BTXRD and FracAtlas using image-level splits and simulated boxes under centered covering and predefined off-center conditions. At $512\times512$, PGA-UNet obtains image-level merged Dice scores of 0.782 and 0.727 under covering prompts, compared with 0.760 and 0.593 for a U-Net baseline with a binary prompt channel. Under the tested off-center condition, its Dice scores are 0.774 and 0.713. These comparisons are descriptive point estimates from the fixed test splits. Ablations show dataset-dependent effects of the conditioning components, while four additional splits characterize variability for PGA-UNet alone. The model contains approximately 2.96 million parameters. The findings support PGA-UNet as a compact design for controlled box-prompted segmentation. Evaluation assumes a correctly localized lesion and does not establish performance with clinician-drawn or off-target prompts.
\end{abstract}

\begin{keywords}
bone X-ray, interactive segmentation, medical image segmentation, prompt-guided learning, Attention U-Net
\end{keywords}

\titlepgskip=-15pt
\maketitle
